\section{Estimación puntual. Insesgadez y mínima varianza}

\begin{ejemplo}
    Sea $(X_1, \ldots, X_n)$ una muestra aleatoria simple de una variable aleatoria $X$ con distribución $\cc{P}(\lm)$, $\lm > 0$. La media muestral $\ol{X}$ es un estimador de $\lm$. Además, cualquier función medible de la muestra que sea independiente del parámetro $\lm$ y tome valores positivos es un estimador del parámetro.
    \begin{equation*}
        \ol{X} = \frac{1}{n} \sum_{i=1}^{n} X_i \geq 0
    \end{equation*}
    donde la positividad de $\ol{X}$ se debe a que $X_i \geq 0$ para todo $i = 1, \ldots, n$.

    Por lo tanto, $\ol{X}$ es un estimador de $\lm$ y cualquier función medible de la muestra que sea independiente del parámetro $\lm$ y tome valores positivos es un estimador del parámetro.
\end{ejemplo}

\begin{ejemplo}
    Sea $(X_1, \ldots, X_n)$ una muestra aleatoria simple de una variable aleatoria $X$ con distribución $\cc{B}(1, p)$, $p \in (0, 1)$. La media muestral $\ol{X}$ es un estimador de $p$. Además, cualquier función medible de la muestra que sea independiente del parámetro $p$ y tome valores en el intervalo $(0, 1)$ es un estimador del parámetro.
    \begin{equation*}
        \ol{X} = \frac{1}{n} \sum_{i=1}^{n} X_i \in \left[0, 1\right]
    \end{equation*}
    donde la pertenencia al intervalo $\left[0, 1\right]$ se debe a que $X_i \in \{0, 1\}$ para todo $i = 1, \ldots, n$. Por lo tanto, $\ol{X}$ es un estimador de $p$ y cualquier función medible de la muestra que sea independiente del parámetro $p$ y tome valores en el intervalo $\left[0, 1\right]$ es un estimador del parámetro.
\end{ejemplo}

\begin{prop}
    Sea $X$ una v.a. con función de distribución en una familia de distribuciones paramétricas, $F \in \{F_\theta : \theta \in \Theta\}$ y $(X_1, \ldots, X_n)$ una m.a.s. de $X$. Sea $T = T(X_1, \ldots, X_n)$ un estimador de $\theta$. El error cuadrático medio (ECM) de $T$ se define como:
    \begin{equation*}
        ECM_T(\theta) := \bb{E}_\theta\left[(T - \theta)^2\right]
    \end{equation*}

    El ECM puede descomponerse en términos de la varianza y una función denominada sesgo:
    \begin{equation*}
        ECM_T(\theta) = Var_\theta(T) + B^2_T(\theta)\qquad \text{donde } B_T(\theta) := \bb{E}_\theta[T] - \theta
    \end{equation*}
    \begin{proof}
        Tenemos que:
        \begin{align*}
            ECM_T(\theta) &= \bb{E}_\theta\left[(T - \theta)^2\right] \\
            &= \bb{E}_\theta\left[(T - \bb{E}_\theta[T] + \bb{E}_\theta[T] - \theta)^2\right] \\
            &= \bb{E}_\theta\left[(T - \bb{E}_\theta[T])^2 + 2(T - \bb{E}_\theta[T])(\bb{E}_\theta[T] - \theta) + (\bb{E}_\theta[T] - \theta)^2\right] \\
            &= Var_\theta(T) + 2(\bb{E}_\theta[T] - \theta)\bb{E}_\theta[T - \bb{E}_\theta[T]] + (\bb{E}_\theta[T] - \theta)^2 \\
            &= Var_\theta(T) + 2(\bb{E}_\theta[T] - \theta)\cancel{(\bb{E}_\theta[T] - \bb{E}_\theta[T])} + (\bb{E}_\theta[T] - \theta)^2 \\
            &= Var_\theta(T) + (\bb{E}_\theta[T] - \theta)^2 \\
            &= Var_\theta(T) + B^2_T(\theta)
        \end{align*}
    \end{proof}
\end{prop}

\begin{ejemplo}
    Sea $(X_1, \ldots, X_n)$ una muestra aleatoria simple de alguna población. Probar que, en general, la cuasidesviación típica no es un estimador insesgado de la desviación típica poblacional.
    \begin{proof}
        Sea $(X_1, \ldots, X_n)$ una muestra aleatoria simple de $X\rightsquigarrow \cc{N}(\mu, \sigma^2)$, $\mu \in \bb{R}$ y $\sigma^2 > 0$. Sabemos que $E[S^2] = \sigma^2$. Queremos ver si $E[S] = \sigma$. Sabemos que:
        \begin{align*}
            Y = \frac{(n-1)S^2}{\sigma^2} \sim \chi^2(n-1)
        \end{align*}
        Por tanto:
        \begin{align*}
            E\left[\sqrt{Y}\right] &= \int_0^\infty \sqrt{y} f_Y(y) dy = \int_0^\infty \sqrt{y}\dfrac{1}{2^{(n-1)/2}\Gamma((n-1)/2)} y^{(n-3)/2} e^{-y/2} dy
            =\\&= \dfrac{1}{2^{(n-1)/2}\Gamma((n-1)/2)} \int_0^\infty y^{n/2-1} e^{-y/2} dy
            \AstIg\\&\AstIg \dfrac{1}{2^{(n-1)/2}\Gamma((n-1)/2)} \cdot \dfrac{\Gamma(n/2)}{(1/2)^{n/2}}
        \end{align*}
        donde en $(\ast)$ hemos aplicado que la función de densidad de la distribución $\chi^2(n)$ integra $1$. Por tanto:
        \begin{align*}
            E[\sqrt{Y}] &= E\left[\sqrt{\frac{(n-1)S^2}{\sigma^2}}\right] = \dfrac{\sqrt{n-1}E[S]}{\sigma}
        \end{align*}

        Por tanto:
        \begin{align*}
            E[S] = \dfrac{\sigma}{\sqrt{n-1}} E[\sqrt{Y}] = \dfrac{\sigma}{\sqrt{n-1}} \cdot \dfrac{1}{2^{(n-1)/2}\Gamma((n-1)/2)} \cdot \dfrac{\Gamma(n/2)}{(1/2)^{n/2}}
            \neq \sigma
        \end{align*}
        Por tanto, $S$ no es un estimador insesgado de $\sigma$.
    \end{proof}
\end{ejemplo}

\begin{ejemplo}
    Sea $(X_1, \ldots, X_n)$ una m.a.s. de una distribución binomial $\cc{B}(1, p)$. Probar que
    \begin{equation*}
        \dfrac{T^2 - T}{n^2 - n}
    \end{equation*}
    con $T = \sum\limits_{i=1}^{n} X_i$ es un estimador insesgado para la función paramétrica $g(p) = p^2$. Probar, además, que para $n = 1$ no existe un estimador insesgado de $p^2$.
    \begin{proof}
        Es directo ver que se trata de un estimador. Veamos que es insesgado, para lo que calculamos su esperanza:
        \begin{align*}
            \bb{E}\left[\dfrac{T^2 - T}{n^2 - n}\right] &= \dfrac{1}{n^2 - n} \bb{E}[T^2 - T] = \dfrac{1}{n^2 - n} \left(\bb{E}[T^2] - \bb{E}[T]\right) =\\&= \dfrac{1}{n^2 - n} \left(Var[T] + \bb{E}[T]^2 - \bb{E}[T]\right) = \dfrac{1}{n^2 - n} \left(np(1-p) + (np)^2 - np\right) =\\&= \dfrac{1}{n^2 - n} \left(n^2p^2 - np^2\right) = p^2
        \end{align*}
        Por tanto, $\dfrac{T^2 - T}{n^2 - n}$ es un estimador insesgado de $p^2$. Veamos ahora que para el valo de $n = 1$ no existe un estimador insesgado de $p^2$. Supongamos que existe un estimador insesgado $T(X)$ de $p^2$ basado en una muestra de tamaño $1$. Entonces, por definición de insesgadez, se cumple que:
        \begin{equation*}
            p^2 = \bb{E}_p[T(X)] = T(0)(1-p) + T(1)p \Longrightarrow p^2 = T(0) - T(0)p + T(1)p
        \end{equation*}
        La expresión de la derecha es una ecuación polinómica de primer grado en $p$, mientras que $p^2$ es un polinomio de segundo grado en $p$. Por tanto, no existe un estimador insesgado de $p^2$ basado en una muestra de tamaño $1$.
    \end{proof}
\end{ejemplo}

\begin{ejemplo}
    Sea $X$ una v.a. con distribución $\cc{P}(\lm)$, ¿existe algún estimador insesgado de $\nicefrac{1}{\lm}$ basado en una muestra de tamaño $1$?\\

    Supongamos que existe un estimador insesgado $T(X)$ de $\nicefrac{1}{\lm}$ basado en una muestra de tamaño $1$. Entonces, por definición de insesgadez, se cumple que:
    \begin{equation*}
        \dfrac{1}{\lm} = \bb{E}_\lm[T(X)] = \sum_{x=0}^{\infty} T(x) \dfrac{e^{-\lm}\lm^x}{x!}
        \Longrightarrow
        e^{\lm} = \sum_{x=0}^{\infty} T(x) \dfrac{\lm^{x+1}}{x!}
    \end{equation*}

    Empleando la serie de Taylor de $e^{\lm}$, obtenemos:
    \begin{equation*}
        \sum_{x=0}^{\infty} \dfrac{\lm^x}{x!} = \sum_{x=0}^{\infty} T(x) \dfrac{\lm^{x+1}}{x!}
    \end{equation*}

    La primera sumatoria tiene un término independiente de $\lm$, mientras que la segunda sumatoria no. Como ha de ser cierta la igualdad para todo $\lm > 0$, se llega a una contradicción. Por tanto, no existe un estimador insesgado de $\nicefrac{1}{\lm}$ basado en una muestra de tamaño $1$.
\end{ejemplo}